\FloatBarrier
\section{Cross-system synthesis: two failures, one decomposition}\label{sec:synthesis}

The two systems were scored separately and will stay separate; what they share is a question and, it turns out, an arithmetic. Both ladders consist of \emph{increment forecasters}: each module forecasts next year's state as this year's state plus a mechanism, and the mechanism is evaluated at whatever the module can substitute for the quantities it needs but does not know. On cod the needed quantity is the catch, which the stock-flow modules are given, and the autonomous modules replace by a constant; on J-17 it is the recharge, which M2 replaces by last year's value, M2m by its climatological mean and the oracle by its realised value. Persistence is the increment forecaster with mechanism zero. The following proposition says exactly how such a forecaster can gain or lose against it.

\subsection{The increment decomposition}\label{sec:decomp}

\begin{proposition}[One-step skill of an increment forecaster]\label{prop:increment}
Let $\Delta_{t+1}=y_{t+1}-y_t$ and let a forecaster issue $\hat y_{t+1}=y_t+m(\hat x_{t+1})$, where $m$ is its mechanism, $x_{t+1}$ the driver at which the mechanism would be evaluated if the driver were known, and $\hat x_{t+1}$ the value substituted at the origin. Write $m_{t+1}=m(x_{t+1})$, $u_{t+1}=\Delta_{t+1}-m_{t+1}$ (the mechanism residual) and $d_{t+1}=m_{t+1}-m(\hat x_{t+1})$ (the substitution error). Then, whenever the second moments exist,
\begin{equation}\label{eq:decomp}
\MSE(\hat y)-\MSE(\pers)\;=\;\underbrace{\E[d_{t+1}^2]-\E[m_{t+1}^2]}_{\text{substitution term}}\;-\;\underbrace{2\,\E\bigl[u_{t+1}\,m(\hat x_{t+1})\bigr]}_{\text{misspecification term}} .
\end{equation}
In particular:
\begin{enumerate}[label=(\roman*),leftmargin=2.2em,itemsep=2pt]
\item \emph{Supplied driver.} If $\hat x_{t+1}=x_{t+1}$ then $d\equiv0$ and $\MSE(\hat y)-\MSE(\pers)=\E[u_{t+1}^2]-\E[\Delta_{t+1}^2]$: the forecaster beats persistence if and only if its mechanism residual is smaller in mean square than the increment itself.
\item \emph{Persisted driver.} If $\hat x_{t+1}=x_t$, the sequence $m_t$ is weakly stationary with mean $\mu$, variance $\sigma^2$ and lag-one autocorrelation $\rho$, and $\E[u_{t+1}m_t]=0$, then
\[\MSE(\hat y)-\MSE(\pers)\;=\;\sigma^2(1-2\rho)-\mu^2,\]
so the forecaster beats persistence if and only if $\rho>\tfrac12-\mu^2/(2\sigma^2)$. For a mean-zero mechanism this is $\rho>\tfrac12$, whatever the mechanism's fidelity.
\item \emph{Climatological driver.} If $m(\hat x_{t+1})=\mu$ and $\E[u_{t+1}]=0$, then $\MSE(\hat y)-\MSE(\pers)=-\mu^2\le0$: the climatological-mechanism forecaster never loses to persistence, and it beats the persisted-driver forecaster of (ii) by $\sigma^2(1-2\rho)$, which is positive if and only if $\rho<\tfrac12$.
\end{enumerate}
\end{proposition}

The proof is in \ref{app:proofs}. The proposition is elementary, and its value is in what it separates. The substitution term depends only on the driver's own predictability from the origin: it is zero when the driver is supplied, it is $\sigma^2(1-2\rho)-\mu^2$ when the driver is persisted, and it is $-\mu^2$ when the driver is replaced by its mean. The misspecification term depends only on how the mechanism's residual co-moves with what the forecaster actually adds; it vanishes when the mechanism is right ($u\equiv0$) or when its residual is unpredictable from the substituted value. A structural forecaster can therefore lose to persistence for two different reasons, and the two systems of this paper lose for one each.

\paragraph{J-17 fails on the substitution term.} The mechanism of M2 is $\beta R+\gamma P$ (the drain term $\delta H_t$ uses a known quantity and is common to the comparison). On the 1934--2023 record the full-sample mechanism $\hat\beta R_t+\hat\gamma P_t$ has lag-one autocorrelation 0.08, standard deviation 9.2\,ft and a mean that is small relative to its spread, so (ii) predicts that M2 loses to persistence by about $\sigma^2(1-2\rho)\approx 71$\,ft$^2$ in mean square, a predicted one-year RMSE near 15.7\,ft against the observed 14.70 (persistence 13.23; the discrepancy is the estimated parameters and the pumpage term, whose own persistence is 0.87). Part (iii) predicts the ordering M2m $\le$ persist $<$ M2, which is what Table~\ref{tab:edwrolling} shows (12.28, 13.23, 14.70\,ft), and it predicts that M2m beats M2 by $\sigma^2(1-2\rho)$, the same 71\,ft$^2$ against an observed $14.70^2-12.28^2=65$\,ft$^2$. The oracle is case (i) with $u$ the map's own residual: $\E[u^2]$ is 7.55$^2$ against $\E[\Delta^2]$ of 13.23$^2$, so the mechanism is a good one, and the only thing wrong with M2 is the value it substitutes for the driver. Fig.~\ref{fig:mechanism}a--b shows the two correlations on which this rests.

\paragraph{Cod fails on the misspecification term.} The stock-flow modules are case (i): the catch is supplied, $d\equiv0$, and their deficit is entirely $\E[u^2]-\E[\Delta^2]$, the mechanism residual against the raw increment. Proposition~\ref{prop:repeller}(ii) says what that residual is in the collapse years: the realised increments $\Delta_{t+1}=S_{t+1}-S_t$ lay 66--161\,kt (coarse regime) and 59--240\,kt (annual landings) below $-C_t$, so $u=\Delta_{t+1}+C_t-g(S_t)<0$ for every nonnegative production curve, and $|u|>|\Delta_{t+1}|$ exactly when $g(S_t)>C_t$, that is, whenever the fitted production exceeds the supplied catch, which is the situation of a module told that removals are falling. The autonomous modules substitute a constant for the catch, which is formally case (ii) with $\rho$ undefined and $\mu$ the training-mean catch; after 1992 that constant exceeds the production of the depleted stock by an order of magnitude, the map is driven to the floor, and $u$ equals the observed biomass at every origin. Neither failure has anything to do with the predictability of a driver: the apparent net production of the assessed series is strongly autocorrelated (lag-one autocorrelation 0.75 on Specification A under the coarse regime, 0.71 under annual landings, 0.91 on Specification B), and the attribution of the cod result to a ``weakly autocorrelated surplus production'' in an earlier combination of these studies is withdrawn. The cod modules fail because the class cannot represent the transitions that carry the error, and the operating characteristics confirm that the test would have seen the structure had it been there.

\subsection{When climatology overtakes persistence}\label{sec:crossover}

The same arithmetic settles the relation between the two baselines, which the five-year J-17 result raised.

\begin{remark}[Horizon crossover of the baselines]\label{rem:crossover}
If $y$ is weakly stationary with variance $\sigma_y^2$ and lag-$h$ autocorrelation $\rho_h$, and the climatological mean is known, then $\MSE_h(\pers)=2\sigma_y^2(1-\rho_h)$ and $\MSE_h(\mathrm{mean})=\sigma_y^2$, so $\RMSE_h(\pers)/\RMSE_h(\mathrm{mean})=\sqrt{2(1-\rho_h)}$ and the training mean overtakes persistence at the first horizon at which $\rho_h<\tfrac12$. With an estimated mean the ratio is shifted by the estimation variance of the mean, which is small relative to $\sigma_y^2$ for long training windows.
\end{remark}

Fig.~\ref{fig:mechanism}c tests the remark on all three series with the baselines recomputed at every horizon $h=1,\dots,5$ on the registered data and on the origin sets of the scored ladders (the $h=1$ and $h=5$ values reproduce Tables~\ref{tab:rolling} and \ref{tab:edwrolling} exactly; SI S23). On J-17 the head's autocorrelation falls from 0.63 at lag one to 0.35 at lag two, so the remark predicts that climatology overtakes persistence between one and two years; the observed ratios are 0.82, 1.07, 1.22, 1.27 and 1.26 against predicted 0.86, 1.14, 1.31, 1.37 and 1.36, the crossover is at $h=2$, and the five-year result of Section~\ref{sec:edwards} is the plateau of that curve rather than an isolated finding. On cod the SSB series is not stationary over the scored windows (a collapse by an order of magnitude and a partial rebuild), the training mean of a collapsing series is a poor estimate of anything, and the remark fails in the direction it should: the observed ratios stay between 0.18 and 0.58 at every horizon while the stationary formula would predict a crossover between four and five years on Specification A. The criterion is therefore a test as well as a prediction. A basin for which $\sqrt{2(1-\hat\rho_h)}$ tracks the observed baseline ratio is one on which a mean-reverting baseline is a defensible multi-year default; a series on which it fails is one on which neither baseline should be trusted at long horizons, and on which the right question is which regime the series is in.

\begin{figure}[!htb]
\centering
\includegraphics[width=\textwidth]{fig7_mechanism}
\caption{The two mechanisms behind the cross-system result. (a) J-17 annual recharge is near-white: $R_t$ against $R_{t-1}$, $r=0.17$. (b) The J-17 head increment tracks same-year recharge: $\Delta H_t$ against $R_t$, $r=0.74$; a balance that persists $R_t$ therefore substitutes a nearly uninformative value for the quantity that explains the increment (Proposition~\ref{prop:increment}(ii)). (c) Ratio of rolling-origin persistence RMSE to training-mean RMSE by horizon on the scored origin sets: J-17 observed (black) against the stationary prediction $\sqrt{2(1-\rho_h)}$ of Remark~\ref{rem:crossover} (grey, dashed), with the crossover at $h=2$; Northern cod Specifications A and B never cross within five years, and the stationary prediction does not apply to them.}
\label{fig:mechanism}
\end{figure}

\subsection{What transfers and what does not}\label{sec:transfer}

Table~\ref{tab:cross} sets the two systems side by side without pooling them. What transfers is the test and the reading of its outcome: a process-based increment forecaster must be scored against persistence and against its own next-simpler rung, on the predictand management uses, at the horizons advice uses, under a protocol stated in advance and disclosed where it was completed late; and when it fails, equation~\eqref{eq:decomp} says where to look, at the driver's predictability from the origin or at the mechanism's residual on the transitions that carry the error. What also transfers is the positive content of each failure. On J-17, the map nowcasts, so the operational question is how to obtain a recharge forecast, and the answer within the information set is that nothing available on 31 December supplies one, which directs effort toward within-year updating rather than toward the balance. On cod, the catch-supplied modules carry a short-memory growth term with one-step skill in a low-biomass regime on one specification, which is a testable claim for a prospective design, and the autonomous modules are refuted with high power, which retires one class of argument about this stock. What does not transfer is any number. The verdicts, margins, autocorrelations and power figures belong to their series, and a reader whose system has a more persistent driver than J-17's recharge, or a predictand that is stationary where cod is not, should expect a different verdict and should re-score rather than inherit one.

\begin{table}[!htb]
\centering\small
\caption{The two systems side by side. Nothing in this table is pooled; each column is one series under its own protocol. ``Failing term'' refers to equation~\eqref{eq:decomp}.}
\label{tab:cross}
\begin{tabularx}{\textwidth}{@{}>{\raggedright\arraybackslash}p{0.21\textwidth}LL@{}}
\toprule
 & Northern cod (Spec.\ A; Spec.\ B) & Edwards J-17 \\
\midrule
Predictand & assessed SSB, kt; retrospective vintage & measured annual-mean head, ft \\
Structural claim & surplus production; catch-driven stock flow; depensation; residual memory & one-pool water balance; residual and delay variants; climate-informed recharge \\
Driver at the origin & catch supplied (M2--M4) or held at training mean (M1, M1b) & recharge and pumpage persisted (M2), climatological (M2m), forecast from climate indices, or realised (oracle) \\
Protocol & C: 5\% band, both horizons, comparator and persistence & E: point rule at $h=1$, persistence and next-simpler causal model \\
One-year persistence RMSE & 98.0; 84.4\,kt & 13.23\,ft \\
Best structural module at $h=1$ & M1b 114.8 (ratio 1.17); M1 119.5 (1.42) & M2m 12.28 (declined by class clause); M1 12.84 (retained) \\
Retained & none; none & AR(1) under E (margin within noise); none under C \\
Failing term & misspecification: collapse transitions outside the class; floor forecasts after 1992 & substitution: $\rho(R)=0.17<\tfrac12$; the map nowcasts at 7.55\,ft \\
Lag-one autocorrelation of the mechanism & apparent net production 0.75; 0.91 & flux mechanism 0.08; recharge 0.17 \\
Baseline crossover (Remark~\ref{rem:crossover}) & none within five years; stationary formula inapplicable & at $h=2$, as predicted \\
What the test could see & power 97--99\% for a true M1; 9--11\% for a true M2; specificity 97--99\% in class & oracle 7.55\,ft, 43\% below persistence; climate indices add $\le0.13$\,ft \\
\bottomrule
\end{tabularx}
\end{table}
